% ── methods/finetune.tex ──────────────────────────────────────────────────────
\subsection{Downstream Fine-Tuning}
\label{sec:finetune}

\paragraph{Shared bottleneck adapter.}
All downstream tasks attach a bottleneck adapter
$A_\psi: \mathbb{R}^d \to \mathbb{R}^d$
(Linear $\to$ GELU $\to$ Linear, bottleneck 128, zero-initialised last layer)
on top of the frozen encoder, adding $<1\%$ extra parameters.

\paragraph{AMP classification.}
A two-layer MLP maps mean-pooled adapter output to a binary AMP logit.
Training data mirrors AMPlify (3{,}338 pos + 3{,}338 neg) for
an apple-to-apple comparison; the encoder is fine-tuned with
differential learning rates ($\eta_\text{head} = 1.5\times10^{-4}$,
$\eta_\text{enc} = 3\times10^{-6}$).

\paragraph{MIC regression.}
GRAMPA provides minimum inhibitory concentrations across 20 bacterial
species~\citep{grampa}.
We evaluate two head types on frozen JEPA embeddings:
(i) \textbf{FiLM-MLP}: a two-layer MLP with FiLM species modulation
$(\hat{y} = \text{MLP}((1{+}\gamma_b)\odot A_\psi(f_\theta(\mathbf{x}))+\beta_b))$;
(ii) \textbf{Transformer}: a two-layer Transformer encoder head
with the same FiLM modulation.
Both use a learned species embedding $e_b \in \mathbb{R}^{64}$ and predict
$\log_2(\text{MIC}/[\mu\text{g\,mL}^{-1}])$.
A formal ESM-2 (35M) FiLM-MLP baseline is trained on the same split with the
same protocol.

\paragraph{QMAP homology-aware regression.}
We evaluate JEPA-AMP on QMAP~\citep{qmap2026} using the official package
(\texttt{qmap-benchmark==0.1.1}), its five predefined splits, and its
log-space metric routine.
Training data for each split is filtered with QMAP's train mask.
We train three variants: (i) frozen embedding + Ridge; (ii) head-only
fine-tune; and (iii) a conditional property head with FiLM modulation over
a shared property-key embedding, enabling one model to predict both
E.~coli MIC and HC50.
All QMAP runs use three random seeds (42, 7, 123); split-level ranges and
paired deltas are archived in \texttt{eval\_results/qmap\_stats\_pack}.
